% NA TOHLE SE RADĚJI ANI NEKOUKEJTE... :-D

\usepackage{xargs}
\usepackage{tikz}
\usetikzlibrary{positioning, math, calc, shadows.blur, trees}

\definecolor{othelloboard}{RGB}{46,90,55}
\definecolor{whitestone}{RGB}{240,240,240}
\definecolor{blackstone}{RGB}{30,30,30}

\tikzset{
    stone/.style={draw=black},
    possible stone/.style={opacity=0.3, font={\footnotesize}},
    shadowed stone/.style={stone, blur shadow={shadow blur steps=7, shadow xshift=0.7pt, shadow yshift=-0.7pt}},
    white stone/.style={fill=whitestone, text=black},
    white shadowed stone/.style={white stone, shadowed stone},
    white possible stone/.style={possible stone, white stone},
    black stone/.style={fill=blackstone, text=white},
    black shadowed stone/.style={black stone, shadowed stone},
    black possible stone/.style={possible stone, black stone}
}

\newcommand{\drawstone}[4]{
    \draw[#3]
        (\numexpr#2-1\relax, \numexpr-#1+1\relax)
        circle (10pt)
        node {#4};
}

\newcommand{\whitestone}[2]{\drawstone{#1}{#2}{white shadowed stone}{}}
\newcommand{\blackstone}[2]{\drawstone{#1}{#2}{black shadowed stone}{}}
\newcommand{\whitepossiblestone}[3]{\drawstone{#1}{#2}{white possible stone}{#3}}
\newcommand{\blackpossiblestone}[3]{\drawstone{#1}{#2}{black possible stone}{#3}}

\newenvironmentx{othellogame}[3][1=show]{
    \begin{tikzpicture}[
            indice/.style={color=black!60, font={\small}},
            scale=0.8
        ]
        \def\rows{#2}
        \def\cols{#3}

        % draw desk
        \draw[shift={(-0.5, 0.5)}, very thick, step=1, fill=othelloboard, rounded corners=0mm] (0, 0) grid (\cols, -\rows) rectangle (0, 0);

        % draw indices for rows
        \ifnum\pdfstrcmp{#1}{show}=0
            \foreach \r in {0,...,\numexpr\rows-1\relax} {
                \pgfmathsetmacro\result{\r+1}
                \node [indice] at (-1, -\r) {\pgfmathprintnumber{\result}};
            }

            % draw indices for cols
            \foreach \c in {0,...,\numexpr\cols-1\relax} {
                \pgfmathsetmacro\result{\c+1}
                \node [indice] at (\c, 1) {\pgfmathprintnumber{\result}};
            }
        \else
            \relax
        \fi
}{\end{tikzpicture}}
